\honest{Mutual Information, and Density in Thingspace}{Eliezer Yudkowsky}{2008}

X has eight equally likely states and Y four, so finding X takes three questions and finding Y two. If X and Y are either both odd or both even, then once you know X, one more question settles Y: four questions in all, not five. The difference is the mutual information, $H(X) + H(Y) - H(X,Y)$, here one bit. \nb{The count assumes that the 16 allowed pairs are equally likely. The post leaves this unsaid; ``Conditional Independence, and Naive Bayes'' later states it.}

Next I give Y the probabilities from my last post, and a new variable Z the probabilities 3/8 and 5/8. Two variables have zero mutual information exactly when every entry of their joint table is the product of the marginals, and that is exactly when learning one tells you nothing about the other. ``Where there is no mutual information, there is no Bayesian evidence, and vice versa.'' The entropy of such a joint table equals the sum of the separate entropies. I print the sum to be computed and leave the calculation to you. \nb{The sum as printed lacks its minus sign. Done correctly, both sides come to 2.704 bits.}

In a second table, Z1 and Y2 occur together more often than their marginals predict, so Y2 is evidence for Z1, and there must be mutual information. I am ``confident, though I haven't actually checked,'' that the joint entropy is smaller than the sum. \nb{It is: 2.664 bits against 2.704, so the mutual information is 0.040 bits.}

Now the words. Why have a word for ``human'' when we could say ``mortal featherless biped''? Because if single properties already have efficient codes, a word for a conjunction is shorter only when the properties are found together more often than chance, and that is exactly when some of them can be inferred from others. \nb{The saving on the pairs that occur more often than chance is paid for by longer codes for the pairs that occur less often. In the post's second table, Z1Y2 saves 0.415 bits, Z1Y3 costs 1.585 bits more, and the average saving is 0.040 bits.} ``Human'' earns its place because a featherless talking biped is likely to be ``poisonable by hemlock, or have broad nails, or be overconfident.'' \nb{Rosch and colleagues had stated the psychological version in 1976: ``The world is structured because real-world attributes do not occur independently of each other.''}

So ``wiggin'' is worth having only if green eyes go with black hair, or if wiggins share other properties. Defining a word implies a promise that it will help inference. If the properties go together no more often than chance, the word ``is a lie''; ``Even if you do not call the word a lie, it is surely an error.''

My conclusion: carve reality at its joints by drawing boundaries ``around concentrations of unusually high probability density in Thingspace.'' \nb{In the next post Yudkowsky says that a qualification was ``deliberately left out'': the boundaries must be simple. That post also calls the rule ``a bit chicken-and-eggish.''}
